\section{Description de l'installation et manipulations}

\subsection{Experience 1, DMA}

Les interets du DMA (voir Ref. \cite{protocoleDSC}:
\\\\
DMA (Dynamic Mechanical Analysis) : mesure des coefficients mécaniques - module de
stockage (élastique) et module de perte (viscoélastique) - lorsque l’échantillon subit une
contrainte périodique.
\\\\

\begin{figure}[ht]
\centering
	\includegraphics[scale=0.6]{schema_pohl.png}
	\captionof{figure}{ Schéma du pendule de Pohl, ref : \cite{schemapohl} }
\label{fig:schemapohl}
\end{figure}

\subsection{Experience 2, DSC}

Les interets du DSC (voir Ref. \cite{protocoleDSC}:
\\\\
"DCS (Differential Scanning Calorimetry) : mesure des flux de chaleurs (exothermique et
endothermique) lorsque l’échantillon subit une rampe de température. La DCS donne accès
aux températures de transition de phase (transition vitreuse, fusion, cristallisation) ainsi qu’aux chaleurs latente".
\\\\
\begin{figure}[ht]
\centering
        \includegraphics[scale=0.6]{schema_pohl.png}
        \captionof{figure}{ Schéma du pendule de Pohl, ref : \cite{schemapohl} }
\label{fig:schemapohl}
\end{figure}

